\section{Confluencia espectral, informacional y termodin\'amica de \(\log3\)}

La identidad

\[
\log3=C_1+C_2-2C_0
=S_{\TRIT}/k_B
=E_P/(k_B\Theta_{\rm clk}),
\]

relaciona tres realizaciones distintas. El primer miembro es un defecto
residual, el segundo una entrop\'ia de decisi\'on y el tercero una raz\'on
entre energ\'ia y temperatura. Su igualdad demuestra la conservaci\'on del
escalar bajo el cambio de realizaci\'on, sin identificar los espacios de
estados.

La primera igualdad se deduce de los tres canales
residuales
\[
Z_r(s)=\sum_{\substack{n\ge1\\n\equiv r\ ({\rm mod}\ 3)}}n^{-s},
\qquad r=0,1,2.
\]
Entonces
\[
Z_0(s)=3^{-s}\zeta(s),\quad
Z_1(s)+Z_2(s)=(1-3^{-s})\zeta(s),\quad
Z_1(s)-Z_2(s)=L(s,\chi_{-3}).
\]
Si \(C_r\) denota la parte finita de \(Z_r\) en \(s=1\), la expansi\'on
de \(3^{-s}\) da
\begin{equation}
C_0+C_1+C_2=\gamma,qquad
-2C_0+C_1+C_2=\log3,qquad
C_1-C_2=\frac{\pi}{3\sqrt3}.
\label{eq:residual-three-covectors}
\end{equation}
Los tres covectores son linealmente independientes. Recuperan, de una sola
descomposici\'on residual, el t\'ermino finito total, la dilataci\'on de la
clase cero y la orientaci\'on entre las dos clases no nulas. \(\log3\) es
por ello un invariante de cambio de escala antes de ser interpretado como
entrop\'ia o temperatura.

\section{Clasificaci\'on epistemol\'ogica de las constantes derivadas}

Las contribuciones del cap\'itulo se clasifican en tres modalidades:

\begin{enumerate}
\item asignaci\'on de un operador interno a una constante conocida, como en
el caso de Catalan;
\item evaluaci\'on compuesta exacta de dos caracteres internos, como
\(L(1,\chi_{12})\);
\item obtenci\'on de una constante a partir de un campo o de una familia
construidos internamente, como \(\gamma_{\mathbb Q(\sqrt3)}\) o \(A_k\).
\end{enumerate}

La procedencia matem\'atica determina el estatuto de recuperaci\'on o
novedad; la proximidad decimal no interviene en dicha clasificaci\'on.

\begin{remark}[Control negativo]
El car\'acter \(\chi_{12}\) debe tener exactamente el patr\'on
\((+,-,-,+)\) en \((1,5,7,11)\); cualquier otro patr\'on invalida
\cref{eq:L1chi12,eq:L2chi12}. La identidad
\cref{eq:odd-zeta-bendersky} y la factorizaci\'on
\cref{eq:dedekind-sqrt3} son controles independientes.
\end{remark}

\subsection*{Alcance lógico y dominio de validez}
\(J^2=-I\), Catalan, \(\chi_{12}\), sus valores \(L\), Euler--Kronecker y la torre Bendersky--Glaisher: \emph{Teorema interno exacto}. La reuni\'on en una ret\'icula \(3/4/12\) es \emph{Formalizaci\'on nueva} que conserva las procedencias matemáticas.

